Todo trabajo de graduación requiere un resumen, pues sirve para que el lector adquiera una visión completa de la materia tratada. Se sugiere colocarlo antes de la introducción porque, en su paso por el índice, el lector ya habrá obtenido una idea general del trabajo, por lo que decidirá si sigue leyendo o no. El resumen aborda una visión general del problema, los objetivos, la metodología y los resultados, sin explicar nada.  

Al final, se incluye una lista de palabras clave (tres mínimo; cinco máximo). Se escriben con minúsculas, separadas por coma. Pueden ser individuales o compuestas, por ejemplo: «educación a distancia», «medicina alternativa», «legislación laboral». Para elegirlas, puede apoyarse de herramientas de IA que ayuden a identificar cuáles serían las mejores palabras clave para motores de búsqueda académicos.

La extensión del resumen es de 250 palabras, por lo que debe ser conciso en la información que brinda.

\palabrasclave{Inteligencia artificial, STEM, ingeniería, \textit{keywords in English}}